%% Generated by verification/make_figures.py from regular-graph-harmonic-quotient-examples.tex.
%% Figure fig:harmonic-quotient_d3reg8a-K4_30213201 (replaces figures/harmonic-quotient_d3reg8a-K4_30213201-G2.tex).
%% Gamma_2 = d3reg8a: 8 vertices, 12 edges (E_2 in the example).
%% phi from the table at line 877. Vertex x is filled with color hq<phi(x)>.
%% Layout: circle, cycle 0,1,...,7.
\begin{tikzpicture}[hqgraph]
\node[hqv, fill=hq3] (v0) at (0.000,1.770) {0};
\node[hqv, fill=hq0] (v1) at (-1.252,1.252) {1};
\node[hqv, fill=hq2] (v2) at (-1.770,0.000) {2};
\node[hqv, fill=hq1] (v3) at (-1.252,-1.252) {3};
\node[hqv, fill=hq3] (v4) at (-0.000,-1.770) {4};
\node[hqv, fill=hq2] (v5) at (1.252,-1.252) {5};
\node[hqv, fill=hq0] (v6) at (1.770,-0.000) {6};
\node[hqv, fill=hq1] (v7) at (1.252,1.252) {7};
\draw[hqe] (v0) -- (v1);
\draw[hqe] (v0) -- (v2);
\draw[hqe] (v0) -- (v7);
\draw[hqe] (v1) -- (v2);
\draw[hqe] (v1) -- (v3);
\draw[hqe] (v2) -- (v3);
\draw[hqe] (v3) -- (v4);
\draw[hqe] (v4) -- (v5);
\draw[hqe] (v4) -- (v6);
\draw[hqe] (v5) -- (v6);
\draw[hqe] (v5) -- (v7);
\draw[hqe] (v6) -- (v7);
\end{tikzpicture}%
